%! AP Statistics — Unit 9, Lesson 9.3 — Warm-Up (Answer Key)
\documentclass[10pt]{article}
\usepackage{apstats-article}
\usepackage{apstats-key}

\begin{document}

\pageheader{Unit 9, Lesson 9.3}{Warm-Up --- Answer Key}
\namedateperiod

\vspace{0.15in}

A researcher used a random sample of 40 adults in a city to study the relationship between
hours of exercise per week ($x$) and resting heart rate (bpm, $y$).
Computer output gave a 95\% confidence interval for the slope: $(-2.8,\ -0.4)$.

\begin{enumerate}

  \item Interpret this confidence interval in context. Use the sentence frame below.

        \textit{``We are 95\% confident that the true slope of the population regression line
        relating \blank{4cm} to \blank{4cm} for \blank{4cm} is between \blank{2cm} and
        \blank{2cm} \blank{3cm}.''}

        \vspace{0.05in}

        \ans{We are 95\% confident that the true slope of the population regression line
        relating resting heart rate to hours of weekly exercise for adults in this city is
        between $-2.8$ and $-0.4$ bpm per hour of exercise per week.}

        \vspace{0.1in}

  \item Does this interval provide evidence that $\beta \ne 0$? Explain briefly.

        \ans{Yes. Since 0 is not contained in the interval $(-2.8,\ -0.4)$, this provides
        evidence that $\beta \ne 0$; there is a negative linear relationship between weekly
        exercise and resting heart rate for adults in this city.}

        \vspace{0.1in}

  \item If the sample size had been 160 instead of 40, how would the \textbf{width} of the
        confidence interval change? Why?

        \ans{The width would decrease. A larger sample size reduces the standard error $SE_b$,
        which makes the margin of error smaller and produces a narrower confidence interval,
        all else remaining the same.}

\end{enumerate}

\end{document}
